% !TeX root = ../main.tex
%%% ---------------------------------------------------------------------------
%%% 方法
%%%
%%% 幻灯片上的公式要能被 3 秒内看懂：
%%%   · 只放论点必需的那一两个式子，推导放备用页
%%%   · 符号用 ../common/acadmath.sty 里的宏（\vect \mat \norm …），
%%%     与论文完全一致，听众看完报告再看论文不需要重新建立对应关系
%%%   · 关键项用 \alert{} 标出来，边讲边指
%%% ---------------------------------------------------------------------------
\section{方法}

\begin{frame}{跨层注意力门控：让每个位置自己决定信任哪一层}
  \begin{block}{门控权重}
    \vspace{-0.6em}
    \begin{equation*}
      \mat{G}_l = \sigma\!\left(
        \mat{W}_2 \relu\!\left(
          \mat{W}_1 \left[ \operatorname{GAP}(\featmap_l) \,\Vert\,
                           \operatorname{GAP}(\tilde{\featmap}_{l+1}) \right]
        \right)
      \right)
    \end{equation*}
  \end{block}

  \vspace{0.5em}
  \begin{block}{融合}
    \vspace{-0.6em}
    \begin{equation*}
      \featmap_l' = \alert{\mat{G}_l} \odot \featmap_l
                  + (\mathbf{1} - \alert{\mat{G}_l}) \odot
                    \mathcal{D}(\tilde{\featmap}_{l+1})
    \end{equation*}
  \end{block}

  \vspace{0.6em}
  \begin{itemize}
    \item $\mat{G}_l \to \mathbf{0}$ 时退化为标准特征金字塔
          \;$\Rightarrow$\; 本方法是它的\alert{严格推广}
    \item 额外参数量仅 \num{0.6}\,M，不改变检测头
  \end{itemize}
\end{frame}

\begin{frame}{门控融合不会放大方差，因此训练是稳定的}
  \begin{theorem}[方差不增]
    设 $\featmap_l$ 与 $\mathcal{D}(\tilde{\featmap}_{l+1})$ 互不相关、
    方差分别为 $\sigma_1^2$ 与 $\sigma_2^2$，则对任意
    $\mat{G}_l \in [0,1]^{H_l \times W_l}$，
    \begin{equation*}
      \Var\!\left[\featmap_l'\right] \le \max\setof{\sigma_1^2, \sigma_2^2}.
    \end{equation*}
  \end{theorem}

  \pause
  \vspace{0.6em}
  \begin{proof}[证明思路]
    逐点展开得 $g^2\sigma_1^2 + (1-g)^2\sigma_2^2$，是 $g$ 的二次函数，
    在 $[0,1]$ 上的最大值只能在端点取得。\qedhere
  \end{proof}

  \vspace{0.4em}
  {\footnotesize\color{AcadMuted}
   完整证明与紧性讨论见备用页与论文附录 A。}

  \note{如果时间紧，这一页只念定理结论，证明思路跳过去，
        有人问再翻备用页。}
\end{frame}

\begin{frame}{可变形对齐消除层级间的感受野偏移}
  \begin{columns}[T, onlytextwidth]
    \begin{column}{0.46\textwidth}
      \begin{equation*}
        \mathcal{D}(\tilde{\featmap})(\vect{p})
          = \sum_{k=1}^{9} w_k\,
            \tilde{\featmap}\!\left(\vect{p} + \vect{p}_k
                                    + \alert{\Delta \vect{p}_k}\right)
      \end{equation*}
      \vspace{0.5em}
      \begin{itemize}
        \setlength{\itemsep}{0.6em}
        \item 偏移量 $\Delta \vect{p}_k$ 由 $3\times3$ 卷积回归
        \item 与门控串联，先对齐再加权
      \end{itemize}
    \end{column}
    \begin{column}{0.50\textwidth}
      \centering
      \placeholder{\linewidth}{4.4cm}{figures/deform-align.pdf}
    \end{column}
  \end{columns}
\end{frame}
